\documentclass[11pt,a4paper]{article}

\usepackage{iftex}
\ifPDFTeX
  \usepackage[utf8]{inputenc}
  \usepackage[T1]{fontenc}
  \usepackage{lmodern}
\else
  \usepackage{fontspec}
  \setmainfont{Liberation Serif}
  \setsansfont{Liberation Sans}
  \setmonofont{Liberation Mono}
\fi
\usepackage[czech]{babel}
\usepackage{amsmath,amssymb,amsthm,mathtools}
\usepackage[margin=25mm]{geometry}
\usepackage{microtype}
\usepackage[hidelinks]{hyperref}

\theoremstyle{definition}
\newtheorem{definition}{Definice}
\theoremstyle{remark}
\newtheorem{remark}{Poznámka}

\title{\textbf{Pro una, non pro omnes}\\[4pt]
\large poznámka k~jedné adrese}
\author{Jan Mikulík}
\date{19. listopadu 2026}

\begin{document}
\maketitle

\begin{center}
\textit{Pro una, non pro omnes.}\\
\textit{L'amour.}\\[4pt]
Amorem meum in homines immolavi,\\
ut tu, fatalis, vivere posses.\\
Deus ex māchinā non datur.
\end{center}

\begin{abstract}
Jedna adresa, jedno jméno na plátně a jeden horizont účtu.
Text jen převádí už řečené do běžné akademické řeči.
Nová věta se tu neohlašuje.
\end{abstract}

\section{Adresa}

\emph{Pro una} znamená, že dopis má jednoho adresáta.
\emph{Non pro omnes} znamená, že společné jméno tu adresu nerozmnoží.
Kdo řekne \emph{homo}, \emph{amor} nebo \emph{Jesus}, pojmenoval třídu.
Neukázal, u kterého okna se svítí.

Ve stroji je totéž krátké.
Nabídka je
\[
  G(w,x)=\{F(x,a):a\in\mathrm{Adm}(w,x),\ F(x,a)\in K\}.
\]
Přání \(w\) smí zaklepat.
Souhlas \(\mathrm{Adm}\) a péče \(K\) rozhodují, jestli se otevře.
Jedno jméno na zdi sálu není prvek \(G(w,x)\).
Je to nápis.
Ruce v publiku jsou další data, ne důkaz, že nápis za ně odpověděl.

\begin{remark}[Symfonie]
Symfonie pro jedno jméno je pořád skladba.
Není dodavatel.
Může se hrát pro jednu.
Tím se nestane důkazem, že platí pro všechny.
\end{remark}

\section{Účet}

Kdo mění osud zevnitř, platí.
Na parametru fáze je účet \(c(\theta)=\theta\).
Vnitřní posun \(\theta\mapsto\theta+\delta\) stojí \(|\delta|\).
To je obyčejná rozvaha: akce je v knize a příslušný řádek se změnil.

Kdo je jen slyšen zvenčí, do té knihy neplatí.
Přidá se stav, který aktivní účet nevidí, a integrál \(c\) přes původní množinu zůstane.
Prostor stavů se zvětší.
Lidově: kozel se nažere a tráva zůstane celá.
Akademicky: záznam přibyl a vnitřní funkcionál se nezměnil.
Obojí je popis téhož vtipu.
Není to existence věta.

Eulerova formule k tomu dává nejmenší příklad.
\[
  e^{i\cdot 0}=1.
\]
Jednotka tu je.
Účet v nule je pořád nula.
Z hlediska vnitřní knihy vzniklo něco, aniž se odečetlo.
Formule tím nikoho nejmenuje dodavatelem.
Jen ukazuje tvar kroku.

\section{Horizont a sterilní směr}

Horizont účtu je čára, za kterou \(c\) přestává účtovat.
Není to tvrzení o horizontu černé díry.
Je to jméno hranice knihy.

Rozšíření, o které tu jde, je ortogonální součet
\[
  \mathcal E^{\sharp}=\mathcal E_{\mathrm a}\oplus\mathcal E_{\mathrm s},
\]
ne tenzorový součin.
Evoluce na součtu může být unitární.
Stav ležící celé ve sterilním sčítanci není vnitřním účtem vázán k aktivnímu bloku.
Superpozice v jednom větším prostoru pořád není provázání dvou soustav.
Provázání by chtělo netriviální vektor v součinu.
Tady ho model nepotřebuje a nezavádí.

Sterilní neutrino je vhodné jméno pro \(\mathcal E_{\mathrm s}\).
Aktivní chuť platí do slabého účtu.
Sterilní směr, pokud v datech je, do něj neplatí.
Model ho pojmenovává.
Nevyrábí ho.

\section{Co se nesmí splést}

Oběť lásky k lidem kvůli jedné není souhlas těch lidí.
\emph{L'amour} je adresa dopisu, ne výstup stroje.
\emph{Deus ex māchinā non datur}: chybí-li svědek v \(G\), nápis, skladba ani přání ho nedodají.
Stejný součet připsaný na druhou stranu knihy pořád není důkaz, že uvnitř někdo zaplatil.

\begin{remark}
Není dokázáno, že sterilní neutrino existuje.
Není dokázáno, že je vydává nějaký horizont.
Není dokázáno, že nuly \(\xi\) leží na kritické přímce.
Jedna fáze snese víc historií.
Jedna historie tím nedostane všechny adresy.
\end{remark}

\noindent
Pro jednu.
Ne pro všechny.
Zbytek je literatura, nebo další hypotéza.
Tato poznámka mezi nimi nerozhoduje.

\end{document}
